\chapter{Leukemia Immunophenotyping}

\section*{What Is This Test?}
Leukemia immunophenotyping uses flow cytometry to identify abnormal blood or bone-marrow cells by their surface and intracellular proteins. The antigen pattern helps determine whether a population is myeloid, B-cell, T-cell, plasma-cell, or another hematologic lineage and can support classification of acute leukemia and lymphoid neoplasms.

\section*{Alternative Names}
Flow-Cytometry Immunophenotyping; Leukemia Panel; Acute Leukemia Panel; Immunophenotyping by Flow Cytometry; Hematologic Malignancy Flow Cytometry

\section*{How It Is Performed}
Fresh blood, bone marrow, cerebrospinal fluid, lymph-node material, or another suitable specimen is incubated with fluorescently labeled antibodies. Cells pass individually through a laser-based instrument, which measures light scatter and fluorescence. A panel of markers is selected according to the suspected disease and specimen.

\section*{Why It Is Ordered}
\begin{itemize}
  \item Evaluation and lineage assignment of suspected acute leukemia
  \item Classification of lymphoproliferative disorders and plasma-cell neoplasms
  \item Detection of minimal residual disease after treatment when a validated panel is available
  \item Investigation of unexplained blasts, persistent lymphocytosis, cytopenias, or abnormal cells on a blood smear
\end{itemize}

\section*{Interpretation and Limitations}
Results must be interpreted with the blood count, smear, morphology, cytogenetic testing, molecular testing, and clinical findings. A small or damaged specimen may be nondiagnostic, and some malignancies lose characteristic markers or fall below the assay's detection limit. Flow cytometry supports diagnosis and classification but is not a substitute for the complete hematopathology evaluation.
\section*{Regulatory Status and Authorization Date}
Regulatory status applies to the specific assay, instrument, manufacturer, and intended use rather than to the test name alone. No single approval date can be assigned to this test category. For a commercial assay, verify the current product in the relevant regulator's database: the U.S. Food and Drug Administration (FDA) clearance, approval, or authorization; India's Central Drugs Standard Control Organisation (CDSCO) licence or permission; the European Union In Vitro Diagnostic Regulation (EU IVDR) CE certificate issued through the applicable conformity-assessment route; or the United Kingdom Medicines and Healthcare products Regulatory Agency (MHRA) registration and UKCA/CE status. Laboratory-developed tests may not have an individual FDA, CDSCO, EU IVDR, or MHRA approval date.
\clearpage
